% Standalone includable proof; packages amsmath and amssymb suffice.
\section*{Independent one-primary complex boundary calculation}

Fix an actual selected zero \(\rho\) of \(g=2\xi\), of its complete
order \(m\geq1\). Set
\[
 \begin{gathered}
 n=m+1,\qquad h(s)=(s-\rho)^m,\qquad t=0,\\
 f(s)=\Phi_h(s)=\frac{(s-\rho)^n-(-\rho)^n}{n},\qquad
 c=-\frac{(-\rho)^n}{n}.
 \end{gathered}
 \tag{SPR.1}
\]
In particular \(f(0)=0\); its constant has not been changed.
Write \(y=s-\rho\), retaining this as a specified coordinate
isomorphism. The original lattice basis is
\(1,s,\ldots,s^{m-1}\), and the translated basis is
\(1,y,\ldots,y^{m-1}\). All coefficient vectors are columns.

For \(0\leq a,b<m\), define
\[
 P_{ab}=
 \begin{cases}\binom ba\rho^{b-a},&a\leq b,\\0,&a>b.\end{cases}
 \qquad
 (P^{-1})_{ab}=
 \begin{cases}\binom ba(-\rho)^{b-a},&a\leq b,\\0,&a>b.\end{cases}
 \tag{SPR.2}
\]
The binomial identity \(s^b=(y+\rho)^b\) proves that \(x_y=P x_s\).
Substitution \(y=s-\rho\) proves the displayed inverse; equivalently,
for \(a\leq b\), their product entry equals
\[
 \binom ba\rho^{b-a}
 \sum_{r=0}^{b-a}\binom{b-a}{r}(-1)^r
 =\begin{cases}1,&a=b,\\0,&a<b.\end{cases}
\]
These are isomorphisms of the complete \(m\)-dimensional coefficient
spaces; no primary jet is discarded.

\subsection*{Exact division and connection}

The polynomial identity, for every \(0\leq b<m\), is
\[
 f(s)s^b=h(s)q_b(s)+r_b(s),\qquad
 q_b(s)=\frac{(s-\rho)s^b}{n},\qquad r_b(s)=cs^b.
 \tag{SPR.3}
\]
Here \(\deg r_b<m\), so this is the unique monic Euclidean division.
The original complex is
\[
 \mathbb C[u,s]\ \xrightarrow{\,L=u\partial_s+h(s)\,}\
 \mathbb C[u,s]\,ds .
 \tag{SPR.4}
\]
Its cokernel is free on the retained monomials, because \(L\) raises
the degree of a nonzero polynomial by exactly \(m\), and monic
subtraction gives a unique remainder of degree less than \(m\).
Subtracting \(Lq_b\) in (SPR.3) gives
\[
 [f(s)s^b]=[cs^b-uq_b'(s)],\qquad
 q_b'(s)=\frac{(b+1)s^b-b\rho s^{b-1}}{n};
 \tag{SPR.5}
\]
the last term is zero when \(b=0\). Its degree is below \(m\).
Consequently the full connection on this lattice, localized at \(u\),
is
\[
 \nabla_u=\partial_u-\frac{cI}{u^2}+\frac{B_s}{u},
 \qquad
 B_s(s^b)=\frac{(b+1)s^b-b\rho s^{b-1}}n.
 \tag{SPR.6}
\]
For completeness the complex connection in degrees \(1,0\) is
\(\partial_u-f/u^2\) and \(\partial_u-f/u^2+1/u\), respectively.
Indeed
\([L,\partial_u-f/u^2]=-\partial_s-h/u=-L/u\).
Thus the degree-one connection composed with \(L\) equals \(L\)
composed with the degree-zero connection. The induced connection is
exactly (SPR.6) by (SPR.5).

The polynomial operator defining \(B_s\) is
\(((s-\rho)\partial_s+1)/n\). Hence
\[
 B_y=P B_sP^{-1}
   =\operatorname{diag}(\beta_0,\ldots,\beta_{m-1}),\qquad
 \beta_a=\frac{a+1}{n}.
 \tag{SPR.7}
\]
This calculation diagonalizes the entire connection by a constant
invertible matrix, since the second-order coefficient is \(cI\).

\subsection*{Ordered contour translation}

Fix a simply connected sector in \(u\ne0\), with its retained branch
\(\log u=\log|u|+i\alpha\). Put
\[
 \theta_j=\frac{\alpha+\pi}{n}+\frac{2\pi j}{n},\qquad
 \ell_j=\{r e^{i\theta_j}:0\leq r<\infty\},\qquad
 \Gamma_j=-\ell_0+\ell_j\quad(1\leq j\leq m).
 \tag{SPR.8}
\]
Every ray is oriented outwards. Let \(\Gamma_j^R\) be its truncation
at radius \(R\), and let \(\rho+\Gamma_j^R\) be its translate,
oriented from \(\rho+Re^{i\theta_0}\) through \(\rho\) to
\(\rho+Re^{i\theta_j}\). Define the endpoint connectors
\[
 K_{j,R}(a)=Re^{i\theta_j}+a\rho,\qquad 0\leq a\leq1.
 \tag{SPR.9}
\]
For every polynomial \(p\), the one-form
\(e^{f(s)/u}p(s)\,ds\) is entire. Cauchy's theorem applied to the
closed path consisting of the original contour, the final connector,
the reverse translated contour, and the reverse initial connector
therefore gives
\[
 \int_{\Gamma_j^R}e^{f/u}p\,ds
 -\int_{\rho+\Gamma_j^R}e^{f/u}p\,ds
 =\int_{K_{0,R}}e^{f/u}p\,ds
   -\int_{K_{j,R}}e^{f/u}p\,ds .
 \tag{SPR.10}
\]
One can equivalently apply Cauchy's theorem to the two ray strips:
their connecting segments from \(0\) to \(\rho\) cancel with their
opposite signs. This proves the identity even if the displayed path
self-intersects.

To justify the limit, on either connector set
\(s=Re^{i\theta_j}+a\rho\). Then
\[
 \frac{(s-\rho)^n}{nu}
 =-\frac{R^n}{n|u|}
   +\frac1{nu}\sum_{r=1}^n
       \binom nr R^{n-r}e^{i(n-r)\theta_j}
                      (a-1)^r\rho^r .
 \tag{SPR.11}
\]
For fixed \(u\), the real part of the last sum is bounded above by
\(K(1+R^{n-1})\), uniformly in \(a\in[0,1]\) and in the finitely
many indices \(j\). If \(b=\deg p\), this gives an upper bound for
the absolute value of either connector integral of the form
\[
 |\rho|\,K_p(1+R)^b
 \exp\left(\operatorname{Re}(c/u)
       -\frac{R^n}{n|u|}+K(1+R^{n-1})\right)
 \ \longrightarrow\ 0 .
 \tag{SPR.12}
\]
When \(\rho=0\), both connectors are constant and their integrals
already vanish. The same estimates, with constants uniform on compact
subsets of the chosen \(u\)-sector, also prove convergence on the
untranslated and translated rays. Thus
\[
 \int_{\Gamma_j}e^{f(s)/u}p(s)\,ds
 =\int_{\rho+\Gamma_j}e^{f(s)/u}p(s)\,ds .
 \tag{SPR.13}
\]
The homotopy \(s\mapsto s+a\rho\) changes the directions at infinity
by a quantity tending to zero, so it preserves each of the labelled
decaying sectors. Formulae (SPR.10)--(SPR.12) prove that the homotopy
respects the convergent periods and their original orientations.
No deformation through a growing end has been used.

\subsection*{The complete period matrix and analytic monodromy}

Put \(\zeta=e^{2\pi i/n}\), and define the \(m\)-by-\(m\) matrix
\[
 W_{ja}=\zeta^{j(a+1)}-1
 \quad(1\leq j\leq m,\ 0\leq a<m),\qquad
 D(u)_{aa}
   =n^{\beta_a-1}e^{\pi i\beta_a}
               \Gamma(\beta_a)u^{\beta_a}.
 \tag{SPR.14}
\]
The branch of \(u^{\beta_a}\) is precisely the one in (SPR.8).
On the ray of angle \(\theta_j\), radial substitution
\(v=r^n/(n|u|)\) proves directly
\[
 \int_{\ell_j} e^{y^n/(nu)}y^a\,dy
 =e^{i(a+1)\theta_j}\frac{(n|u|)^{\beta_a}}n
       \int_0^\infty e^{-v}v^{\beta_a-1}\,dv
 =\zeta^{j(a+1)}D(u)_{aa}.
 \tag{SPR.15}
\]
The integral defining \(\Gamma(\beta_a)\) is positive and finite,
since \(0<\beta_a<1\). Translation (SPR.13), followed by the full
binomial expansion of \(s^b=(y+\rho)^b\), therefore gives the exact
period matrix in the original monomial lattice:
\[
 (\Pi_s)_{jb}=\int_{\Gamma_j}e^{f(s)/u}s^b\,ds,\qquad
 \boxed{\ \Pi_s(u)=e^{c/u}W D(u)P.\ }
 \tag{SPR.16}
\]
To prove \(W\) invertible, suppose \(Wa=0\) and set
\(F(z)=\sum_{r=1}^m a_{r-1}(z^r-1)\). This polynomial vanishes
at all the nontrivial \(n\)th roots of unity by \(Wa=0\), and at
\(1\) by its definition. Its degree is at most \(n-1\), so \(F=0\);
the coefficient of \(z^r\) then gives \(a_{r-1}=0\) for every \(r\).
All diagonal entries of \(D(u)\) and the scalar exponential are
nonzero. Thus (SPR.16) is an invertible matrix for every \(u\ne0\)
on the branch. Its determinant retains the original constant:
\[
 \det\Pi_s
 =(\det W)n^{-m/2}e^{\pi i m/2}
     \left(\prod_{a=0}^{m-1}\Gamma(\beta_a)\right)
        u^{m/2}e^{mc/u},
 \tag{SPR.17}
\]
because \(\det P=1\) and \(\sum_a\beta_a=m/2\).

Analytic continuation once counterclockwise around \(u=0\) multiplies
each \(u^{\beta_a}\) by \(\zeta^{a+1}\) and leaves \(e^{c/u}\)
unchanged. Consequently
\[
 \Pi_s^{\mathrm{cont}}=M_{\mathrm{per}}\Pi_s,\qquad
 M_{\mathrm{per}}
   =W\operatorname{diag}(\zeta,\zeta^2,\ldots,\zeta^m)W^{-1}.
 \tag{SPR.18}
\]
This is also the ordered-contour transformation. During continuation
\(\theta_j\) increases by \(2\pi/n\), so
\[
 \Gamma_j\longmapsto
 \begin{cases}
    \Gamma_{j+1}-\Gamma_1,&1\leq j<m,\\
    -\Gamma_1,&j=m.
 \end{cases}
 \tag{SPR.19}
\]
The matrix with these rows sends \(W_{ja}\) to
\(\zeta^{a+1}W_{ja}\), proving that (SPR.18) and (SPR.19)
use the same orientation and continuation convention.

The spectrum in (SPR.18) consists of every nontrivial \(n\)th root,
each once. The matrix is diagonalizable, has exact order \(n\), and
has no invariant vector. In particular
\[
 \ker(M_{\mathrm{per}}-I)=0,\qquad
 \det(zI-M_{\mathrm{per}})
       =\frac{z^n-1}{z-1}.
 \tag{SPR.20}
\]
This assertion concerns this one factor with \(t=0\). It does not
assert the absence of invariant vectors in a tensor power.

Here are the horizontal-connection conventions explicitly.
Writing \(K_s=-cI/u^2+B_s/u\), differentiation of (SPR.16) gives
\(\Pi_s'=\Pi_sK_s\). Therefore \(F_s=\Pi_s^{-1}\) satisfies
\(F_s'=-K_sF_s\), so its columns are a fundamental set of
horizontal sections of \(\nabla_u=\partial_u+K_s\).
Continuation gives
\[
 F_s^{\mathrm{cont}}=F_sM_{\mathrm{per}}^{-1}.
 \tag{SPR.21}
\]
Equivalently, the fundamental matrix
\[
 \widehat F_s
   =P^{-1}\operatorname{diag}
       (e^{-c/u}u^{-\beta_0},\ldots,e^{-c/u}u^{-\beta_{m-1}})
 \tag{SPR.22}
\]
has horizontal monodromy
\(\operatorname{diag}(\zeta^{-1},\ldots,\zeta^{-m})\).
Both conventions have no eigenvalue \(1\). The inverse in (SPR.21)
must be retained when passing from periods to horizontal sections.

\subsection*{Ramification, the scalar irregular factor, and Stokes data}

The constant coordinate change (SPR.2) gives an exact direct sum
of scalar connections
\[
 \nabla_{y,a}=\partial_u-\frac c{u^2}+\frac{\beta_a}{u},
 \qquad 0\leq a<m.
 \tag{SPR.23}
\]
This is an actual meromorphic direct sum, valid on the punctured
disc, not merely a formal asymptotic one. Pull it back under the
specified cover \(u=v^n\). The result is
\[
 \nabla_{v,a}
    =\partial_v-\frac{nc}{v^{n+1}}+\frac{a+1}{v}.
 \tag{SPR.24}
\]
The single-valued meromorphic gauge
\(x_a=v^{-(a+1)}z_a\) cancels the displayed integer residue and
leaves
\[
 \partial_v-\frac{nc}{v^{n+1}}
 \tag{SPR.25}
\]
on every summand. Its horizontal solutions are
\(z_a=e^{-c/v^n}\) times constants. The scalar irregular
exponential has been retained in (SPR.23)--(SPR.25);
\(e^{c/u}\) has not been treated as a meromorphic gauge at zero.
In particular, if \(c\ne0\), these formulas do not declare the
original connection regular singular.

There are no nonidentity Stokes transition matrices: the displayed
constant diagonalization, followed on the cover by the displayed
meromorphic diagonal gauge, already provides the same exact model
on every sector. Its comparison with that model is the identity,
and comparisons on overlaps are therefore the identity as well.
There is only one scalar irregular exponential, so no distinct
exponential factors require different sectorial splittings.
The possible branch transitions of \(u^{-\beta_a}\) are exactly
the finite monodromy already calculated in (SPR.21); they are not
additional Stokes transitions.

After the cover, the analytic monodromy is \(I\), because
\(M_{\mathrm{per}}^n=I\). Its unipotent logarithm is therefore
exactly zero:
\[
 N_{\mathrm{boundary}}=\log I=0 .
 \tag{SPR.26}
\]
On the original punctured disc the horizontal local system has
no invariant sections under a positive loop, by (SPR.21).
This is a statement about the complex analytic local system
of the computed connection. It has not identified that local
system with finite-field inertia or with a Frobenius representation.

\subsection*{Original arithmetic jets and their exact connection}

In \(E_h=\mathbb C[s]/(s-\rho)^m\), let \(A_s\) be multiplication
by \(s\). Let \(N\) be multiplication by \(y\) on
\(\mathbb C[y]/y^m\), in the complete ordered basis:
\[
 N_{a+1,a}=1\quad(0\leq a<m-1),\qquad
 \text{all other entries }0.
 \tag{SPR.27}
\]
The coordinate isomorphism (SPR.2) proves
\[
 P A_sP^{-1}=\rho I+N,\qquad
 N_s:=A_s-\rho I=P^{-1}NP .
 \tag{SPR.28}
\]
Here \(N^m=0\), and \(N^{m-1}\ne0\) for every \(m\geq1\),
with \(N^0=I\) when \(m=1\). In particular \(N=0\) precisely
when \(m=1\). The complete nilpotent arithmetic action is retained.

For \(a<m-1\), the difference of the two successive diagonal entries
of \(B_y\) is \(1/n\). Thus on every basis vector
\[
 [B_y,N]=\frac Nn,\qquad
 \boxed{\ [B_s,A_s]=\frac{A_s-\rho I}{n}.\ }
 \tag{SPR.29}
\]
For the top vector both sides of the first identity are zero.
This proves the identity on the whole primary algebra, including its
top jet. It constructs an exact relationship between the arithmetic
nilpotent and the computed boundary residue:
\(N_s=n[B_s,A_s]\).
The vanishing logarithm (SPR.26) has consequently not erased the
nonzero arithmetic nilpotent when \(m>1\).

Finally retain the entire original source multiplier
\(v_h(s)=g(s)/(s-\rho)^m\). Since \(m\) is the exact zero order,
the full primary Taylor unit and its matrix are
\[
 \begin{gathered}
 v_h(\rho+y)\bmod y^m
   =\sum_{r=0}^{m-1}v_r y^r,\qquad
 v_r=\frac{g^{(m+r)}(\rho)}{(m+r)!},\qquad v_0\ne0,\\
 U_y=\sum_{r=0}^{m-1}v_rN^r,\qquad U_s=P^{-1}U_yP.
 \end{gathered}
 \tag{SPR.30}
\]
Writing \(R=v_0^{-1}(U_y-v_0I)\), the exact inverse is
\[
 U_y^{-1}=v_0^{-1}\sum_{j=0}^{m-1}(-R)^j,
 \tag{SPR.31}
\]
because \(R^m=0\). No Taylor coefficient permitted by the full
primary quotient has been removed.

The original source map
\(\mathcal T_h(p)=p(D)F_h\) has Mellin transform \(p\,v_h\);
therefore its full jet is \(U_s x_s\), not merely \(x_s\).
In the source-weighted jet coordinates \(q_s=U_sx_s\),
the exact multiplication and connection coefficients are
\[
 A_s^\theta=U_sA_sU_s^{-1}=A_s,\qquad
 B_s^\theta=U_sB_sU_s^{-1},\qquad
 C_s^\theta=cI .
 \tag{SPR.32}
\]
The equality for \(A_s^\theta\) holds because \(U_s\) is a polynomial
in \(A_s-\rho I\). In translated weighted coordinates,
(SPR.29) implies, by induction on \(r\),
\([B_y,N^r]=(r/n)N^r\). Hence the full correction is
\[
 B_y^\theta
   =B_y-\frac1n
        \left(\sum_{r=1}^{m-1}r v_rN^r\right)U_y^{-1},
 \qquad
 [B_y^\theta,\rho I+N]=\frac Nn .
 \tag{SPR.33}
\]
The correction is a polynomial in \(N\), so the last commutator
follows on all jets. The period matrix on the same weighted jet
coordinates is precisely \(\Pi_sU_s^{-1}\); the original source
and its unit have thus been attached by an explicit invertible map.
Here \(U_s\) extends the marked-fibre jet map to the free coefficient
lattice by the displayed constant matrix in its retained basis.
This statement does not identify that extension with multiplication
by \(v_h(s)\) on the twisted de Rham quotient away from \(u=0\).
In fact the polynomial Taylor representative
\(v_{\mathrm{pol}}(s)=\sum_{r=0}^{m-1}v_r(s-\rho)^r\) satisfies
the exact relation
\[
 v_{\mathrm{pol}} Lq
   =L(v_{\mathrm{pol}}q)-u v_{\mathrm{pol}}'q
 \qquad(q\in\mathbb C[u,s]).
 \tag{SPR.34}
\]
This follows by expanding both derivatives using the product rule.
It retains the derivative correction when multiplication is applied
to polynomial representatives of the original complex.
More explicitly, for each fixed \(u\) let
\[
 H_u=\mathbb C[s]\,ds/L(\mathbb C[s])\,ds,\qquad
 \iota_u:E_h\longrightarrow H_u,\quad
 [p]_h\longmapsto[(\operatorname{rem}_h p)\,ds]_L,\qquad
 T_{\upsilon,u}=\iota_u M_{[v_h]_h}\iota_u^{-1}.
 \tag{SPR.35}
\]
The unique remainder argument following (SPR.4) proves that
\(\iota_u\) is a linear isomorphism; its definition uses the
specified degree-less-than-\(m\) representative, so it is independent
of the representative of \([p]_h\). This constructs the exact
coefficient transport \(T_{\upsilon,u}\) on \(H_u\), whose matrix is
\(U_s\). If \(\deg p<m\), perform the ordinary monic division
\(v_{\mathrm{pol}}p=hq+r\), with \(\deg r<m\).
The quotient, when nonzero, has degree at most \(m-2\).
Writing \(\operatorname{red}_{L,u}\) for the unique polynomial
representative of degree less than \(m\), one obtains
\[
 \operatorname{red}_{L,u}(v_{\mathrm{pol}}p)
     =r-uq',\qquad
 T_{\upsilon,u}[p\,ds]_L=[r\,ds]_L .
 \tag{SPR.36}
\]
Indeed \(hq=Lq-uq'\), and \(q'\) already has degree less than
\(m\). For \(m=1\), both factors are constant and \(q=0\).
Thus (SPR.36) calculates the full difference between literal
polynomial multiplication and the coordinate transport on the
retained lattice; it vanishes at the marked fibre \(u=0\).
All calculations above concern the single selected primary factor.
They make no inference from this one-factor invariant calculation
to tensor invariants or to the Riemann hypothesis.
